<%= heading2 %>{Tables}

<%- bexample %>
\begin{table}
  \caption{Simple Table}
  \label{tab:simple}
  \centering
  \begin{tabular}{ll}
    \toprule
    Heading1 & Heading2 \\
    \midrule
    One      & Two      \\
    Three    & Four     \\
    \bottomrule
  \end{tabular}
\end{table}
<%- eexample %>

A table caption goes above the table (\verb|\caption| before the \verb|tabular|), a figure caption below the figure.

Rows and cells can be colored with \verb|\rowcolor| and \verb|\cellcolor| from the \href{https://ctan.org/pkg/colortbl}{colortbl} package, which \href{https://ctan.org/pkg/xcolor}{xcolor} loads with its option \verb|table|.

<%- bexample %>
\begin{table}
  \caption{Table with colored cells}
  \label{tab:colored}
  \centering
  \begin{tabular}{ll}
    \toprule
    \rowcolor{gray!20}
    Test & Result                   \\
    \midrule
    A    & \cellcolor{green!20}pass \\
    B    & \cellcolor{red!20}fail   \\
    \bottomrule
  \end{tabular}
\end{table}
<%- eexample %>

A \verb|\footnote| inside a floating table is lost.
Use \verb|\tablefootnote| from the \href{https://ctan.org/pkg/tablefootnote}{tablefootnote} package instead: it is set as a normal footnote at the bottom of the page.
If the notes should appear directly below the table instead (common in journals), use the \href{https://ctan.org/pkg/threeparttable}{threeparttable} package.

<%- bexample %>
\begin{table}
  \caption{Table with a footnote}
  \label{tab:footnote}
  \centering
  \begin{tabular}{ll}
    \toprule
    Heading1 & Heading2\tablefootnote{A footnote inside a table.} \\
    \midrule
    One      & Two                                                \\
    \bottomrule
  \end{tabular}
\end{table}
<%- eexample %>

A table wider than the text runs into the margin.
\verb|\begin{adjustbox}{max width=\textwidth}| from the \href{https://ctan.org/pkg/adjustbox}{adjustbox} package scales it down to the text width, and only if it is too wide.
Scaling shrinks the font, so for a table with long text in its cells, prefer wrapping the text with the column type \verb|X| of the \href{https://ctan.org/pkg/tabularx}{tabularx} package, and for a very wide table, a landscape page (see \texttt{pdflscape}).

<%- bexample %>
\begin{table}
  \caption{Wide table scaled to the text width}
  \label{tab:wide}
  \centering
<%# [impl->req~wide-table-scaled~1] -%>
  \begin{adjustbox}{max width=\textwidth}
    \begin{tabular}{llllllllllll}
      \toprule
      Jan  & Feb  & Mar  & Apr  & May  & Jun  & Jul  & Aug   & Sep  & Oct   & Nov   & Dec   \\
      \midrule
      1.11 & 2.22 & 3.33 & 4.44 & 5.55 & 6.66 & 7.77 & 8.88  & 9.99 & 10.10 & 11.11 & 12.12 \\
      Rain & Snow & Wind & Sun  & Fog  & Hail & Heat & Storm & Mist & Frost & Sleet & Calm  \\
      \bottomrule
    \end{tabular}
  \end{adjustbox}
\end{table}
<%- eexample %>

<%- include('diagbox.example.en.tex', this); %>
